\section{RL / Agent Infra：当 Actor 变大、Rollout 变长、环境变复杂}

\subsection{RL 的瓶颈越来越像推理问题}
多位嘉宾明确指出，今天 RL infra 的主要卡点是 rollout，而不是训练本身。这一点很关键，因为它解释了为什么推理引擎、缓存管理、环境接口和数据回流，会在 RL 时代重新占据系统设计中心。尤其是模型变大、rollout 变长后，传统 RL 框架面对的是和过去完全不同的 Actor。

\begin{importantbox}{今天的大模型 RL 和传统 RL 的本质差别}
嘉宾拿早年的 MuJoCo / Atari 系统作比较，指出差异主要体现在三点：
\begin{itemize}
  \item Actor 从几兆参数网络，变成了巨大的 LM / VLM；
  \item 环境从相对固定的小型 simulator，变成了更动态、更异构的真实工具和服务；
  \item rollout 不再是廉价 forward，而是一次完整的推理与交互过程。
\end{itemize}
\end{importantbox}

\subsection{长 rollout 和大模型把系统压力成倍放大}
嘉宾们提到了几个非常具体的挑战：模型更大时如何做 RL、rollout 特别长时如何做 RL、multi-modality 和 Agent workload 如何做 RL。这些问题看似算法问题，实际上都直接映射到系统压力。比如 rollout 变长意味着缓存更久、调度更复杂、失败恢复成本更高；环境更复杂意味着训练框架必须支持更异构的数据流和状态流。

\begin{knowledgebox}{为什么 RL infra 需要训推结合}
如果训练框架和推理框架完全割裂，就会出现几个问题：
\begin{itemize}
  \item rollout 结果回传链路过长；
  \item 状态管理分散在多个系统中；
  \item 调试故障时很难定位是在训练、推理还是环境侧；
  \item 算法侧和系统侧往往会 ``大眼瞪小眼''，谁都觉得瓶颈在对方。
\end{itemize}
\end{knowledgebox}

\subsection{Agent RL 把环境构造问题推到台前}
传统 RL 常常假设环境已经给定，但在 Agent 任务中，环境本身就是系统设计的一部分。panel 中多次提到 environment construction、workflow construction、real world task、failure handling 等问题，这说明 Agent RL 不只是 ``在既有环境中学策略''，而是经常要先把环境搭出来，再让模型去学。

\begin{warningbox}{环境如果构造不好，RL 就会学偏}
嘉宾提到一个核心痛点：很多系统在单环境里能训练得很好，一旦跨环境就不泛化。也就是说，环境搭建差异、任务结构差异、奖励形式差异都会直接成为泛化问题。这不是简单增加数据量就能抹掉的误差。
\end{warningbox}

\subsection{本章小结}
RL / Agent Infra 的复杂度，来自三个变量同时变化：Actor 变大、rollout 变长、环境变复杂。于是系统设计的重点也从 ``把训练跑起来'' 转向 ``让训练、推理、环境和数据回流形成可复现、可扩展的闭环''。

